%%%%%%%%%%  scan idx 44  |  folio 31  %%%%%%%%%%
La
formule (1.22) admet les généralisations suivantes~: soit $K$ un $\lambda$-anneau quelconque et par ailleurs arbitraire~; pour $x\in K$, on pose
\begin{equation}\tag{1.23}
\gamma^n(x)=(-1)^n\sum_{i=0}^n(-1)^i\lambda^i(x+n)=\lambda^n(x+n-1)
\end{equation}
(l'identité des deux derniers membres résulte d'une récurrence immédiate). On a en particulier $\gamma^0(x)=1$ et $\gamma^1(x)=x$~; on pose ensuite~:
\begin{equation}\tag{1.24}
\gamma_t(x)=\sum_{n\geq0}\gamma^n(x)t^n\in K[[t]]
\end{equation}
d'où l'on déduit la formule~:
\begin{equation}\tag{1.25}
\gamma_t(x)=\lambda_{t/(1-t)}(x),
\end{equation}
ce qui équivaut à
\begin{equation}\tag{1.25 bis}
\lambda_s(x)=\gamma_{s/(1+s)}(x).
\end{equation}
Ceci prouve que les $\lambda^i$ s'expriment à l'aide des $\gamma^i$ et que $\gamma_t(x+y)=\gamma_t(x)\gamma_t(y)$, et par suite que les applications $\gamma^i$ définissent une nouvelle structure de $\lambda$-anneau sur $K$. La formule (1.22) montre alors que pour $K=\widetilde A$ et tout $x\in\widetilde A$, on a
\begin{equation}\tag{1.22 bis}
\gamma^n(x-\varepsilon(x))=[0,1+(-1)^{n-1}(n-1)!x^n+\cdots]
\end{equation}
et par suite, la \emph{composante} $x^n$ \emph{s'obtient, au facteur} $(-1)^{n-1}(n-1)!$ \emph{près, en réduisant modulo} $\widetilde A^{n+1}$ \emph{l'élément} $\gamma^n(x-\varepsilon(x))$ \emph{de} $\widetilde A^n$.

\paragraph{L'homomorphisme de Chern~:}
Pour tout anneau gradué $A$ vérifiant les conditions précédentes, on va définir un homomorphisme d'anneaux augnentés~:
\begin{equation}\tag{1.26}
\ch:\widetilde A\relbar\joinrel\relbar\joinrel\relbar \widehat{A\otimes\Q}.
\end{equation}
Cet homomorphisme est défini sans ambiguité par la condition d'être un homomorphisme additif, de transformer $1$ en $1$, d'être tel que les
